\documentclass{article} % Define la clase del documento, en este caso, un artículo
\usepackage[letterpaper,margin=3cm]{geometry} % Configura el tamaño del papel y los márgenes del documento
\usepackage{graphicx} % Permite la inserción de imágenes
\usepackage[spanish]{babel}% Activar esta configuración para informes en español, ajusta el idioma del documento
\usepackage[usenames]{color} % Permite el uso de colores definidos por nombre en el documento
\usepackage{hyperref} % Habilita enlaces y referencias dentro del documento
\hypersetup{colorlinks=true, linkcolor = black, citecolor= black} % Configura el color de los enlaces y citas
\usepackage{booktabs} % Proporciona comandos para crear tablas de alta calidad
\usepackage{natbib} % Permite el uso de citas y referencias bibliográficas con diferentes estilos
\usepackage{tikz} % Permite la creación de gráficos y diagramas vectoriales directamente en LaTeX
\usepackage{float} % Para controlar la posición de los elementos flotantes, como imágenes, con la opción [H]
\usepackage{diagbox} % Permite crear celdas con líneas diagonales en tablas
\usepackage{listings} % Permite la inclusión y formateo de código fuente en el documento
\usepackage{xcolor} % Paquete para definir y usar colores en el documento
\usepackage{parskip} % Añade espacio entre párrafos en lugar de sangrías
\usepackage{fancyhdr} % Permite personalizar encabezados y pies de página
\usepackage{amsmath} % Proporciona una amplia variedad de entornos y comandos matemáticos

\pagestyle{fancy} % Usa el estilo fancyhdr
\fancyhf{} % Borra todos los encabezados y pies de página
\renewcommand{\headrulewidth}{0pt}
\renewcommand{\footrulewidth}{0pt} % Desactiva la línea horizontal predeterminada en el pie
\setlength{\headheight}{2cm} % Ajusta la altura del encabezado para hacer espacio para la línea
\fancyhead[L]{\raisebox{0.20cm}{\textbf{Métodos y Técnicas de Construcción}}} % Añade el texto en la parte izquierda del encabezado, subiéndolo ligeramente
\fancyhead[R]{\raisebox{0.1cm}{\includegraphics[width=0.25\linewidth]{LOGO_UNIVERSIDAD.jpg}}} % Añade la imagen en la parte derecha del encabezado y súbela un poco
\fancyhead[C]{\rule{\textwidth}{0.6pt}} % Añade una línea horizontal superior centrada
\fancyfoot[C]{\rule{\textwidth}{0.6pt}} % Añade una línea horizontal en el pie de página centrada
\fancyfoot[R]{\raisebox{-1.5\baselineskip}{\thepage}} % Coloca el número de página a la derecha, con suficiente espacio debajo de la línea
\geometry{top=3cm, bottom=2.5cm} % Ajusta los márgenes superior e inferior

% Definición de colores al estilo Visual Studio Code
\definecolor{codegreen}{rgb}{0.25,0.49,0.48} % Comentarios
\definecolor{codegray}{rgb}{0.5,0.5,0.5} % Números y anotaciones
\definecolor{codepurple}{rgb}{0.58,0,0.82} % Palabras clave
\definecolor{backcolour}{rgb}{0.95,0.95,0.92} % Color de fondo

% Configuración del estilo de las celdas de código
\lstset{
    backgroundcolor=\color{backcolour},   % color de fondo; necesita que el paquete color o xcolor esté cargado
    commentstyle=\color{codegreen},       % estilo de comentarios
    keywordstyle=\color{codepurple},      % estilo de palabras clave
    numberstyle=\tiny\color{codegray},    % estilo de los números de línea
    stringstyle=\color{red},              % estilo de las cadenas de texto
    basicstyle=\ttfamily\small,           % estilo del texto básico
    breakatwhitespace=false,              % ajustes de líneas sólo en espacios en blanco
    breaklines=true,                      % ajustar las líneas si son muy largas
    captionpos=b,                         % posición de la leyenda (abajo)
    keepspaces=true,                      % preserva los espacios en el texto; útil si se usa monoespaciado
    numbers=left,                         % dónde poner los números de línea
    numbersep=5pt,                        % qué tan lejos están los números de línea del código
    showspaces=false,                     % mostrar espacios con subrayados particulares; reemplaza 'showstringspaces'
    showstringspaces=false,               % subrayar los espacios dentro de las cadenas solo
    showtabs=false,                       % mostrar tabulaciones en el código con subrayados particulares
    tabsize=2,                            % tamaños de tabulación a 2 espacios
    language=TeX,                         % lenguaje del código
    morecomment=[l]\#,                    % reconocer # como inicio de comentario en Python
    frame=single,                         % agregar un marco simple alrededor del código
    rulecolor=\color{black}               % color del marco
}

\begin{document}
%----------------------------------------------------------------------------------------
%   PORTADA
%Modificar desde aqui en adelante
%----------------------------------------------------------------------------------------
\begin{titlepage}%Inicio de la carátula, solo modificar los datos necesarios
\newcommand{\HRule}{\rule{\linewidth}{0.5mm}} 
\center 
%----------------------------------------------------------------------------------------
%	ENCABEZADO
%----------------------------------------------------------------------------------------
\includegraphics[width=10cm]{LOGO_UNIVERSIDAD.jpg}\\ % Si esta plantilla se copio correctamente, va a llevar la imagen del logo de la facultad.OBS: Es necesario incluir el paquete: graphicx
\vspace{3cm}
%----------------------------------------------------------------------------------------
%	SECCION DEL TITULO
%----------------------------------------------------------------------------------------
\HRule \\[0.4cm]
{ \huge \bfseries Tarea 1}\\[0.4cm] % Titulo del documento
{ \huge \bfseries Métodos y Técnicas de Construcción}\\[0.4cm] % Titulo del documento
\HRule \\[1.5cm]
 \vspace{5cm}
%----------------------------------------------------------------------------------------
%	SECCION DEL AUTOR
%----------------------------------------------------------------------------------------
\begin{flushright}
    { \textbf{Profesor:}\\
    Jose Tramon\\
    \vspace{0.2cm}
    \textbf{Alumnos:}\\
    Bernardo Caprile Canala-Echevarría\\
    Pedro Tomás Valenzuela Bejares\\
    \vspace{0.2cm}

}
\end{flushright}
\vspace{1cm}
%----------------------------------------------------------------------------------------
%	SECCION DE LA FECHA
%----------------------------------------------------------------------------------------
{\large \textbf{\today}}\\[2cm] % El comando \today coloca la fecha del dia, y esto se actualiza con cada compilacion, en caso de querer tener una fecha estatica, reemplazar el \today por la fecha deseada
\end{titlepage}
%----------------------------------------------------------------------------------------
%  INDICE
%----------------------------------------------------------------------------------------
\newpage
\tableofcontents
\thispagestyle{plain} % Deshabilita el encabezado en la página del índice
\thispagestyle{empty} % Deshabilita el número de página en la página del índice
\newpage

%Se puede agregar un indice de figuras si es nesesario
%\newpage
%\listoffigures 
%\thispagestyle{plain} % Deshabilita el encabezado en la página del índice %
%\thispagestyle{empty}
%\newpage
%----------------------------------------------------------------------------------------
%   ACÁ EMPIEZA EL INFORME
\setcounter{page}{1} % Reinicia el contador de páginas
%-------
\section{Introducción}

En la gestión de proyectos de construcción, es muy importante planificar bien para asegurar que todo se haga a tiempo y de acuerdo a lo previsto. Un paso clave en esta planificación es calcular correctamente los volúmenes de trabajo, lo que se conoce como cubicación. Esto nos permite estimar cuánto tiempo tomará completar cada parte del proyecto. En este caso, se nos ha pedido realizar la cubicación y planificación de las obras que forman parte de una licitación pública, disponible en el portal de Mercado Público bajo el ID 976-10-O123. El proyecto consiste en construir un paso bajo nivel en una nueva vía llamada Diego Barros Ortiz, e incluye trabajos como excavaciones, construcción de muros de contención y cajones de hormigón armado.

La cubicación es esencial para poder crear un cronograma de trabajo usando el software MS Project. Este cronograma deberá mostrar el orden en que se realizarán las actividades, cuánto tiempo tomará cada una, y cuál es la ruta crítica, es decir, las tareas que determinarán el tiempo total del proyecto.
\newpage

\section{Criterios de duración de Actividades}

\subsection{Vialidad}

\begin{table}[h!]
    \centering
    \begin{tabular}{|c|l|c|r|}
    \hline
    \textbf{Item} & \textbf{Descripción} & \textbf{Unidad} & \textbf{Cantidad} \\ \hline
    2.1 & \textbf{Vialidad} & & \\ \hline
    2.1.1 & Remoción de pavimento de hormigón & m² & 361.6 \\ \hline
    2.1.2 & Remoción de pavimento asfáltico & m² & 452 \\ \hline
    2.1.3 & Excavación en terreno de cualquier naturaleza & m³ & 31,904.18728 \\ \hline
    2.1.4 & Formación y construcción de terraplenes & m³ & 40,671 \\ \hline
    2.1.5 & Excavación de escarpe & m³ & 0 \\ \hline
    2.1.6 & Sub base granular CBR 40\% & m³ & 1,475.244 \\ \hline
    2.1.7 & Carpeta granular de rodadura & m³ & 0 \\ \hline
    2.1.8 & Base granular CBR 80\% & m³ & 760.35 \\ \hline
    2.1.9 & Imprimación & m² & 4,620 \\ \hline
    2.1.10 & Riego de liga & m² & 4,445 \\ \hline
    2.1.11 & Concreto asfáltico, capa de rodadura & m³ & 222.25 \\ \hline
    2.1.12 & Concreto asfáltico, capa intermedia (binder) & m³ & 266.7 \\ \hline
    2.1.13 & Aceras de hormigón & m² & 780 \\ \hline
    2.1.14 & Soleras tipo A & m & 1,270 \\ \hline
    2.1.15 & Señal vertical lateral tipo 2 & un & 27 \\ \hline
    2.1.16 & Señal vertical lateral tipo 3 & un & 1 \\ \hline
    2.1.17 & Demarcación, línea de eje segmentada & m & 412 \\ \hline
    2.1.18 & Demarcación, línea lateral continua & m & 230 \\ \hline
    2.1.19 & Demarcación, líneas, achurados, símbolos y leyendas & m² & 11 \\ \hline
    \end{tabular}
    \caption{Cubicación de Vialidad}
\end{table}

En cada partida, dentro de las cuadrillas consideradas para los distintos trabajos, se incluyó tanto al personal de obra como a los operarios encargados de la operación de la maquinaria necesaria. Además, no se consideraron la excavación de escarpe ni la carpeta granular de rodadura, ya que estas no se utilizaban en esta etapa del proyecto.

\subsubsection{Remoción de Pavimento de Hormigón (m²)}

Para esta partida se consideró un rendimiento de 50 $\frac{m^2}{hora}$ por cuadrilla de 3 a 5 personas, dependiendo del equipo y condiciones del terreno. Con 2 cuadrillas se estima un tiempo total de 1 día.

\subsubsection{Remoción de Pavimento Asfáltico (m²)}

Para esta partida se consideró un rendimiento de 70 $\frac{m^2}{hora}$ por cuadrilla de 3 a 5 personas, dependiendo de las herramientas y maquinaria empleada. Con 2 cuadrillas se estima un tiempo total de 1 día.

\subsubsection{Excavación en Terreno de Cualquier Naturaleza (TCN) (m³)}

Para esta partida se consideró un rendimiento de 22 $\frac{m^3}{hora}$ por cuadrilla de 4 a 6 personas. Con 5 cuadrillas se estima un tiempo total de 37 días.

\subsubsection{Formación y Construcción de Terraplenes (m³)}

Para esta partida se consideró un rendimiento de 90 $\frac{m^3}{hora}$ por cuadrilla de 4 a 6 personas. Con 2 cuadrillas se estima un tiempo total de 29 días.

\subsubsection{Sub Base Granular CBR 40\% (m³)}

Para esta partida se consideró un rendimiento de 15.625 $\frac{m^3}{hora}$ por cuadrilla de 4 a 6 personas. Con 2 cuadrillas se estima un tiempo total de 6 días.

\subsubsection{Base Granular CBR 80\% (m³)}

Para esta partida se consideró un rendimiento de 12.5 $\frac{m^3}{hora}$ por cuadrilla de 4 a 6 personas. Con 2 cuadrillas se estima un tiempo total de 4 días.

\subsubsection{Imprimación (m²)}

Para esta partida se consideró un rendimiento de 500 $\frac{m^2}{hora}$ por cuadrilla de 3 a 4 personas. Con 1 cuadrilla se estima un tiempo total de 2 días.

\subsubsection{Riego de Liga (m²)}

Para esta partida se consideró un rendimiento de 600 $\frac{m^2}{hora}$ por cuadrilla de 3 a 4 personas. Con 1 cuadrilla se estima un tiempo total de 1 día.

\subsubsection{Concreto Asfáltico, Capa de Rodadura (m³)}

Para esta partida se consideró un rendimiento de 30 $\frac{m^3}{hora}$ por cuadrilla de 6 a 8 personas. Con 1 cuadrilla se estima un tiempo total de 1 día.

\subsubsection{Concreto Asfáltico, Capa Intermedia (Binder) (m³)}

Para esta partida se consideró un rendimiento de 30 $\frac{m^3}{hora}$ por cuadrilla de 6 a 8 personas. Con 1 cuadrilla se estima un tiempo total de 2 días.

\subsubsection{Aceras de Hormigón (m²)}

Para esta partida se consideró un rendimiento de 40 $\frac{m^2}{hora}$ por cuadrilla de 4 a 6 personas. Con 2 cuadrillas se estima un tiempo total de 2 días.

\subsubsection{Solera Tipo A (m)}

Para esta partida se consideró un rendimiento de 20 $\frac{m}{hora}$ por cuadrilla de 4 a 6 personas. Con 3 cuadrillas se estima un tiempo total de 3 días.

\subsubsection{Señal Vertical Lateral Tipo 2 (UN)}

Para esta partida se consideró un rendimiento de 15 $\frac{UN}{hora}$ por cuadrilla de 3 a 4 personas. Con 1 cuadrilla se estima un tiempo total de 1 día.

\subsubsection{Señal Vertical Lateral Tipo 3 (UN)}

Para esta partida se consideró un rendimiento de 10 $\frac{UN}{hora}$ por cuadrilla de 3 a 4 personas. Con 1 cuadrilla se estima un tiempo total de 1 día.

\subsubsection{Demarcación, Línea De Eje Segmentada (m)}
Para esta partida se consideró un rendimiento de 700 $\frac{m}{hora}$ por cuadrilla de 3 a 4 personas. Con 1 cuadrilla se estima un tiempo total de 1 día.


\subsubsection{Demarcación, Línea Lateral Continua (m)}

Para esta partida se consideró un rendimiento de 700 $\frac{m}{hora}$ por cuadrilla de 3 a 4 personas. Con 1 cuadrilla se estima un tiempo total de 1 día.

\subsubsection{Demarcación, Líneas, Achurados, Símbolos y Leyendas (m²)}

Para esta partida se consideró un rendimiento de 50 $\frac{m^2}{hora}$ por cuadrilla de 3 a 4 personas. Con 1 cuadrilla se estima un tiempo total de 1 día.

\subsection{Estructural}

\begin{table}[h!]
    \centering
    \begin{tabular}{|c|l|c|r|}
    \hline
    \textbf{Item} & \textbf{Descripción} & \textbf{Unidad} & \textbf{Cantidad} \\ \hline
    2.2 & \textbf{ESTRUCTURAS} & & \\ \hline
    2.2.1 & Excavación a Maquina en Puentes y Estructuras & m³ & 22,322.19 \\ \hline
    2.2.2 & Relleno Estructural Permeable & m³ & 50,336.286 \\ \hline
    2.2.3 & Impermeabilización de Estribo y Muros & m² & 3,759.1268 \\ \hline
    2.2.4 & Hormigón G-5 & m³ & 292.746828 \\ \hline
    2.2.5 & Hormigón G-30/SM & m³ & 4,751.637 \\ \hline
    2.2.6 & Acero para Armaduras A63-42H & KG & 651,338.153 \\ \hline
    2.2.7 & Moldajes para Infraestructura & m² & 14,621.523 \\ \hline
    2.2.8 & Losa de Acceso & m³ & 98.34 \\ \hline
    2.2.9 & Barbacanas de Desagüe & UN & 70.4 \\ \hline
    \end{tabular}
    \caption{Cubicación de las Estructuras}
\end{table}

Es importante mencionar, que la losa de acceso y las barbacanas estan consideradas en la partida de hormigón G-30/SM, Hormigón G5, moldajes y Acero para Armaduras A63-42H, por lo que no se consideran en esta tabla.

\subsubsection{Excavación a máquina en Puentes y Estrucutras y Relleno Estructural}
Para estas dos partidas se consideró un rendimiento de 100 $\frac{m^3}{dia}$ por retroexcavadora, consiguiendo 28 y 34 días de trabajo respectivamente.

\subsubsection{Moldajes}
Para esta partida se consideró un rendimiento de 240 $\frac{m^2}{dia}$ por cuadrilla, consiguiendo 31 días de trabajo.

\subsubsection{Acero para armaduras A63-42H}
Para esta partida se consideró un rendimiento de 1.12 $\frac{ton}{dia}$ por cuadrilla, para este trabajo se consideraron 3 cuadrillas, consiguiendo 194 días de trabajo.

\subsubsection{Hormigón G-5}
Para esta partida se consideró un rendimiento de 28 $\frac{m^3}{hora}$, esto de la base que un camión mixer de $7m^3$ se demora aproximadamente 30 minutos en vaciarse. Por la gemometría del proyecto pueden estar vaciandose 2 caminoes al mismo tiempo, esto conlleva que en 1 hora se logren vaciar 4 camiones mixer. Finalmente, se consideraron 2 días de trabajo.

\subsubsection{Hormigón G-30/SM}
Para esta partida se consideró el mismo rendimiento que la partida anterior, por lo que esta tiene un tiempo aproximado de 22 días.

\subsubsection{Impermeabilización de Estribo y Muro}
Para esta partida se consideró un rendimiento de 30 $\frac{m^2}{dia}$ por trabajador, para optimizar el tiempo se consideraron 2 cuadrillas de 3 personas consiguiendo 21 días de trabajo.

\subsubsection{Relleno Estructural Permeable}
Para esta partida se consideró un rendimiento de 100 $\frac{m^3}{dia}$ por retroexcavadora, consiguiendo 34 días de trabajo.



\newpage
\section{Criterios de relación}

Primero, se deben realizar las actividades de preparación del área de trabajo, que incluyen la remoción de pavimento de hormigón y la remoción de pavimento asfáltico. Estas tareas deben completarse antes de proceder con cualquier trabajo de excavación, ya que permiten despejar la superficie de trabajo. Una vez que se ha limpiado el área, se procede con el movimiento de tierras de la parte vial y estructural. La excavación en terreno de cualquier naturaleza es la primera etapa, seguida por la formación y construcción de terraplenes, la cual se realiza para estabilizar el terreno donde se realizarán las siguientes capas.

Después de completar el movimiento de tierras, se pueden comenzar las actividades de estructuras. La excavación a máquina en puentes y estructuras se realiza para preparar las áreas donde se colocarán las fundaciones de los cajones y muros de contención. Posteriormente, se procede con la colocación de los moldajes para la posterior colocación del hormigón G-5. Tras esto, se empieza a colocar la enfierradura y se puede verter el hormigón G-30/SM. Luego, se impermealizan los estribos y muros de contención. Una vez sacado el moldaje se procede a las fases de pavimentación.

Simultaneamente con la etapa estrucutral, se pueden colocar las capas granulares. Primero, se instala la sub base granular CBR 40\%, luego la carpeta granular de rodadura CBR 60\%, y finalmente la base granular CBR 80\%. Con las capas granulares completadas, se pasa a la fase de pavimentación, que incluye la imprimación y el riego de liga, que preparan el suelo para recibir las capas de pavimento. Después, se colocan el concreto asfáltico, capa de rodadura y la capa intermedia. Una vez que el pavimento está instalado, se pueden construir las aceras de hormigón y la solera tipo A.

Finalmente, se colocan el relleno y la losa de acceso. Posteriormente, se puede proceder a las actividades de señalética y demarcación, las cuales incluyen la instalación de señales verticales y la demarcación de las líneas de eje y laterales en el pavimento.

\newpage
\section{Conclusión}

En conclusión, la cubicación de la calle de rodaje del aeropuerto AMB ha sido un proceso fundamental para asegurar una planificación eficiente de los recursos y tiempos requeridos en el desarrollo de la obra. Gracias a este análisis, hemos calculado con los volúmenes de trabajo para las distintas etapas agregando un factor de seguridad del 10 \%, desde la remoción de pavimentos existentes hasta la construcción de terraplenes, capas granulares, pavimentos y elementos estructurales.

Estos cálculos no solo permiten prever las cantidades de materiales necesarios, como hormigón, asfalto y elementos de señalización, sino que también facilitan la asignación adecuada de cuadrillas y maquinaria, optimizando así la ejecución de las actividades. Además, este enfoque garantiza que el proyecto pueda ajustarse a los plazos establecidos, reduciendo riesgos y posibles sobrecostos.

Finalmente, este proceso de cubicación es clave para la creación de una línea de tiempo de obra realista y bien estructurado, que permita organizar las actividades de manera eficiente y evitar demoras. A través de la cubicación, se puede identificar la ruta crítica del proyecto, asegurando que cada tarea se realice en el tiempo adecuado y que los recursos estén disponibles en cada fase. Esto contribuye directamente al éxito del proyecto, permitiendo cumplir con los plazos, optimizar los costos y garantizar que se respeten los estándares de calidad y las normativas exigidas.

\newpage
\bibliographystyle{agsm}
\begin{thebibliography}{}

    \bibitem[ONDAC, 2024a]{ondac_article_161909}
    ONDAC. (2024a). "Portal ONDAC." [Online]. Available at: \url{https://portal.ondac.com/601/w3-article-161909.html} [Accessed 29 September 2024].
    
    \bibitem[ONDAC, 2024b]{ondac_article_165862}
    ONDAC. (2024b). "Portal ONDAC." [Online]. Available at: \url{https://portal.ondac.com/601/w3-article-165862.html} [Accessed 29 September 2024].
    
    \bibitem[ONDAC, 2024c]{ondac_article_61489}
    ONDAC. (2024c). "Portal ONDAC." [Online]. Available at: \url{https://portal.ondac.com/601/w3-article-61489.html} [Accessed 29 September 2024].
    
    \bibitem[ONDAC, 2024d]{ondac_article_61490}
    ONDAC. (2024d). "Portal ONDAC." [Online]. Available at: \url{https://portal.ondac.com/601/w3-article-61490.html} [Accessed 29 September 2024].
    
    \bibitem[ONDAC, 2024e]{ondac_article_161792}
    ONDAC. (2024e). "Portal ONDAC." [Online]. Available at: \url{https://portal.ondac.com/601/w3-article-161792.html} [Accessed 29 September 2024].
    
    \bibitem[ONDAC, 2024f]{ondac_article_169382}
    ONDAC. (2024f). "Portal ONDAC." [Online]. Available at: \url{https://portal.ondac.com/601/w3-article-169382.html} [Accessed 29 September 2024].
    
    \bibitem[ONDAC, 2024g]{ondac_article_161589}
    ONDAC. (2024g). "Portal ONDAC." [Online]. Available at: \url{https://portal.ondac.com/601/w3-article-161589.html} [Accessed 29 September 2024].
    
    \bibitem[Aridos Madesal, 2024]{aridos_base_cbr}
    Aridos Madesal. (2024). "Portal ONDAC." [Online]. Available at: \url{https://aridosmadesal.cl/producto/base-estabilizada-cbr-80-r/} [Accessed 29 September 2024].
    
    \bibitem[ONDAC, 2024h]{ondac_article_161659}
    ONDAC. (2024h). "Portal ONDAC." [Online]. Available at: \url{https://portal.ondac.com/601/w3-article-161659.html} [Accessed 29 September 2024].
    
    \bibitem[ONDAC, 2024i]{ondac_article_64794}
    ONDAC. (2024i). "Portal ONDAC." [Online]. Available at: \url{https://portal.ondac.com/601/w3-article-64794.html} [Accessed 29 September 2024].
    
    \bibitem[ONDAC, 2024j]{ondac_article_161588}
    ONDAC. (2024j). "Portal ONDAC." [Online]. Available at: \url{https://portal.ondac.com/601/w3-article-161588.html} [Accessed 29 September 2024].
    
    \bibitem[ONDAC, 2024k]{ondac_article_92694}
    ONDAC. (2024k). "Portal ONDAC." [Online]. Available at: \url{https://portal.ondac.com/601/w3-article-92694.html} [Accessed 29 September 2024].
    
    \bibitem[ONDAC, 2024l]{ondac_article_64795}
    ONDAC. (2024l). "Portal ONDAC." [Online]. Available at: \url{https://portal.ondac.com/601/w3-article-64795.html} [Accessed 29 September 2024].
    
    \bibitem[ONDAC, 2024m]{ondac_article_161791}
    ONDAC. (2024m). "Portal ONDAC." [Online]. Available at: \url{https://portal.ondac.com/601/w3-article-161791.html} [Accessed 29 September 2024].
    
    \bibitem[ONDAC, 2024n]{ondac_article_161590}
    ONDAC. (2024n). "Portal ONDAC." [Online]. Available at: \url{https://portal.ondac.com/601/w3-article-161590.html} [Accessed 29 September 2024].
    
    \bibitem[ONDAC, 2024o]{ondac_article_172694}
    ONDAC. (2024o). "Portal ONDAC." [Online]. Available at: \url{https://portal.ondac.com/601/w3-article-172694.html} [Accessed 29 September 2024].
    
    \bibitem[Cesmec, 2010]{cesmec_2010}
    CESMEC. (2010). "Portal ONDAC." [Online]. Available at: \url{chrome-extension://efaidnbmnnnibpcajpcglclefindmkaj/https://www.cesmec.cl/medios/DIC/Documentos/VOL_5_DIC_2010.pdf} [Accessed 29 September 2024].
    
    \bibitem[Chile Cubica, 2024]{chilecubica}
    Chile Cubica. (2024). "Portal Chile Cubica." [Online]. Available at: \url{https://www.chilecubica.com/estudio-costos/rendimientos/#:~:text=%C2%BFCu%C3%A1ntos%20m3%20hace%20una%20excavadora,aproximado%20de%205%20m3%2Fjornada.} [Accessed 29 September 2024].
    

\end{thebibliography}

\end{document}